\cotefile{64}
\scanpage{145}
classes d'isomorphie d'objets \unsure{d'icelles|d'icelle} est en corr. biunivoque canonique avec l'ens des applications holomorphes
\makeatletter\tagsleft@true\makeatother
\[
\tau:S\longrightarrow\mathbb C\setminus\mathbb R . \tag*{(5)}
\]
On dit qu'une structure de Torelli est ``directe'' ou ``stricte'' si \deleted{\unsure{l'est|l'ori.}} pour $\unsure{\forall|\mathrm{tt}}\,s\in S$, la base \deleted{de} $z_1,z_2$ de $\unsure{H|X|\mathcal H}_s\otimes_{\mathbb Z}\mathbb R$ \deleted{\unsure{correspond à|correspond|correspond.}}\inserted{définit} l'\deleted{\unsure{orient.|struc.}}\inserted{orientation} de cet espace correspondant à sa structure complexe, i.e si (5) est un morphisme
\[
\tau:S\longrightarrow\mathcal D^+=\{z\in\mathbb C\mid\unsure{J|\mathrm I|\Im}z>0\}. \tag*{(6)}
\]
Donc $\mathcal D^+$ apparaît comme l'espace analytique modulaire pour la structure ``courbe elliptique relative avec str. de Torelli stricte sur espace analytique variable'', la famille modulaire étant donnée par
\[
X_{\mathcal D^+}=(\mathcal D^+\times\mathbb C)/\mathbb Z^2 \tag*{(7)}
\]
où $\mathbb Z^2$ opère sur $\mathcal D^+\times\mathbb C$ par
\[
\begin{aligned}
\deleted{\unsure{m,n}}\;T_{m,n}(\tau,z)&=(\tau,z+m+n\tau)\\
&\hspace{5em}\tau\in\mathcal D^+,\ z\in\mathbb C
\end{aligned}\tag*{(8)}
\]
de sorte que l'on a
\[
\begin{aligned}
(X_{D^+})_\tau &\simeq \mathbb C/(\mathbb Z+\mathbb Z.\tau)\simeq\mathbb C^*/\text{sous-groupe}\\
&\hspace{10em}\text{engendré}\\
&\hspace{10em}\text{par }\deleted{\unsure{q}}\;k=\exp(2i\pi\tau)\\
&\hspace{10em}\text{(\unsure{satisfaisant|on trouve bien} }|k|<1\text{)}
\end{aligned}\tag*{(8)}
\]
% archive sheet number: 144
